\documentclass[11pt,a4paper]{article}
\usepackage{fontspec,amsmath,amssymb,amsthm,microtype}
\setmainfont{lmroman10-regular.otf}[BoldFont=lmroman10-bold.otf,ItalicFont=lmroman10-italic.otf,BoldItalicFont=lmroman10-bolditalic.otf]
\usepackage{xcolor}
\usepackage[most]{tcolorbox}
\usepackage{mathrsfs}
\providecommand{\lt}{<}
\providecommand{\gt}{>}
\setmonofont{DejaVuSansMono.ttf}[BoldFont=DejaVuSansMono-Bold.ttf,ItalicFont=DejaVuSansMono-Oblique.ttf,BoldItalicFont=DejaVuSansMono-BoldOblique.ttf,Scale=0.86]
\newfontfamily\leanfallback{STIXTwoMath-Regular.otf}[Scale=0.86]
\definecolor{leanrule}{gray}{0.62}
\definecolor{leanback}{gray}{0.962}
\definecolor{leantitle}{gray}{0.22}
\definecolor{leanmuted}{gray}{0.38}
\newtcolorbox{leanbox}[1]{enhanced,breakable,colback=leanback,colframe=leanrule,
  boxrule=0.4pt,arc=0.6mm,left=2.4mm,right=2.4mm,top=1.8mm,bottom=1.8mm,
  before skip=0.8em,after skip=1em,fonttitle=\small\bfseries,coltitle=leantitle,
  colbacktitle=leanback,titlerule=0.3pt,title={#1}}
\newlength{\leancharwidth}
\newenvironment{leancode}{\par\smallskip\begingroup\ttfamily\footnotesize\raggedright
  \settowidth{\leancharwidth}{M}\setlength{\parskip}{0pt}\setlength{\parindent}{0pt}
  \frenchspacing}{\par\endgroup\smallskip}
\newcommand{\leanline}[2]{\hangindent=\dimexpr#1\leancharwidth+4\leancharwidth\relax
  \hangafter=1\noindent\hspace*{#1\leancharwidth}#2\par}
\newcommand{\leansource}[1]{\par{\footnotesize\color{leanmuted}#1}\par}
\newcommand{\leanlabel}[1]{\par\smallskip\noindent{\small\bfseries #1}\par}
\usepackage[margin=25mm]{geometry}
\usepackage[unicode,hidelinks]{hyperref}
\hypersetup{pdfauthor={LangevinLean Project}}
\setlength{\emergencystretch}{3em}
\setlength{\parskip}{0.4em}
\setlength{\parindent}{0pt}
\allowdisplaybreaks[2]
\raggedbottom
\title{\texorpdfstring{The KIPLMC1 parameter root-mean-square error bound}{The KIPLMC1 parameter root-mean-square error bound}}
\author{LangevinLean Project}
\date{}
\begin{document}\maketitle
Draft mathematical exposition; these notes may contain mistakes.


The KIPLMC1 Gaussian Markov laws are identified with the shared constructed KLMC laws at the changed clock. The shared kinetic theorem constructs its Brownian reference, proves exact Gibbs stationarity, bounds the local error, and takes the infimum over initial couplings. The particle geometry and sharp parameter concentration then give the paper’s three explicit error terms.



The constructed shared theorem applies to the sum potential and correlated-Gaussian chain. Exact clock cancellation and sharp parameter concentration yield the three explicit physical-time KIPLMC1 error contributions.



Earlier results used here: \cite[Section~3.114]{LangevinLean2026Kinetic}; \cite[Section~3.57]{LangevinLean2026Kinetic}.

\section{\texorpdfstring{The particle noise source with the Gaussian Markov law}{The particle noise source with the Gaussian Markov law}}


The particle coordinate set is \(\mathrm{Fin}(d_\theta)\sqcup(\mathrm{Fin}(N)\times\mathrm{Fin}(d_x))\). The shared kinetic construction now works directly on this arbitrary finite set. Its complete constructed innovation stream has the proved product law of independent spatial blocks, each with the exact correlated two-component Gaussian distribution. The argument instantiates the product-source Markov-law theorem with this Brownian source, its whole-stream law, and an independent standard Gaussian initial velocity. This identifies the kinetic position law constructed above with the position marginal of the Gaussian Markov iterate. The particle clock-rescaling theorem then gives the exact identity between the physical-time KIPLMC1 output and this same specified output for \(V_N\) at friction \(\sqrt N\gamma\) and step \(h/\sqrt N\). The position target is unchanged.



The finite particle coordinate set has cardinality \(d_\theta+Nd_x\), even when the parameter and latent dimensions differ. The generic shared construction applies directly to that set. Its Brownian response law therefore matches the particle Gaussian transition without an artificial dimensional identification.

\begin{leanbox}{Lean: \texttt{canonicalKLMC\_\allowbreak{}eq\_\allowbreak{}gaussianPositionLaw}, \texttt{kiplmc1GaussianPositionLaws\_\allowbreak{}eq\_\allowbreak{}canonical}}
\leansource{\texttt{canonicalKLMC\_\allowbreak{}eq\_\allowbreak{}gaussianPositionLaw} \quad KIPLMC/\allowbreak{}CanonicalExponential.lean, lines 15–26}
\begin{leancode}
\leanline{0}{theorem\ canonicalKLMC\_eq\_gaussianPositionLaw\ \{N\ dθ\ dx\ :\ ℕ\}}
\leanline{4}{(f\ :\ State\ N\ dθ\ dx→ℝ)\ (M\ :\ ℝ≥0)\ (hLip\ :\ LipschitzWith\ M\ (gradient\ f))}
\leanline{4}{(γ\ :\ ℝ)\ (h\ :\ ℝ≥0)\ (ν\ :\ Measure\ (State\ N\ dθ\ dx))\ [IsProbabilityMeasure\ ν]\ (n\ :\ ℕ)\ :}
\leanline{4}{GenericCanonicalCoupling.klmcLaw\ f\ γ\ h\ ν\ n\ =}
\leanline{6}{(ULA.markovLaws\ (klmcStep\ (gradient\ f)\ γ\ h)\ (exponentialNoise\ N\ dθ\ dx\ γ\ h)}
\leanline{8}{((stdGaussian\ (State\ N\ dθ\ dx)).prod\ ν)\ n).map\ Prod.snd}
\end{leancode}
\leansource{\texttt{kiplmc1GaussianPositionLaws\_\allowbreak{}eq\_\allowbreak{}canonical} \quad KIPLMC/\allowbreak{}CanonicalExponential.lean, lines 29–39}
\begin{leancode}
\leanline{0}{theorem\ kiplmc1GaussianPositionLaws\_eq\_canonical\ \{N\ dθ\ dx\ :\ ℕ\}\ (hN\ :\ 0<N)}
\leanline{4}{(U\ :\ IPLA.Joint\ dθ\ dx→ℝ)\ (L\ :\ ℝ≥0)}
\leanline{4}{(hdiff\ :\ Differentiable\ ℝ\ U)\ (hLip\ :\ LipschitzWith\ L\ (gradient\ U))}
\leanline{4}{(γ\ :\ ℝ)\ (h\ :\ ℝ≥0)\ (hγ\ :\ 0≤γ)}
\leanline{4}{(ν\ :\ Measure\ (State\ N\ dθ\ dx))\ [IsProbabilityMeasure\ ν]\ (n\ :\ ℕ)\ :}
\leanline{4}{(kiplmc1GaussianLaws\ U\ γ\ h\ ((velocityTarget\ N\ dθ\ dx).prod\ ν)\ n).map\ Prod.snd\ =}
\leanline{6}{GenericCanonicalCoupling.klmcLaw\ (@IPLA.potential\ N\ dθ\ dx\ U)\ (Real.sqrt\ N*γ)}
\leanline{8}{⟨(h:ℝ)/Real.sqrt\ N,div\_nonneg\ h.coe\_nonneg\ (Real.sqrt\_nonneg\ \_)⟩\ ν\ n}
\end{leancode}
\leanlabel{How the Lean proof follows the argument}
\texttt{kiplmc1\allowbreak{}Step} is the deterministic particle step. \texttt{kiplmc1\allowbreak{}Step\_\allowbreak{}parameter\_\allowbreak{}position} and \texttt{kiplmc1\allowbreak{}Step\_\allowbreak{}particle\_\allowbreak{}velocity} identify its coordinate blocks with the displayed physical equations. \texttt{kiplmc1\allowbreak{}Gaussian\allowbreak{}Position\allowbreak{}Laws\_\allowbreak{}eq\_\allowbreak{}canonical} transfers the complete Gaussian Markov law, including initialization and correlated innovations, to the canonical kinetic sampler at the changed clock.
\end{leanbox}


Earlier results used here: \cite[Section~3.33.6]{LangevinLean2026ULASC}; \cite[Section~3.8.17]{LangevinLean2026ULASC}; \cite[Section~3.1.1]{LangevinLean2026IPLA}; \cite[Section~3.1.4]{LangevinLean2026IPLA}; \cite[Section~3.114.3]{LangevinLean2026Kinetic}; \cite[Section~3.12.2]{LangevinLean2026Kinetic}; \cite[Section~3.36.1]{LangevinLean2026Kinetic}; \cite[Section~3.36.2]{LangevinLean2026Kinetic}; \cite[Section~3.57.14]{LangevinLean2026Kinetic}.

\section{\texorpdfstring{The transformed assumptions and all particle-count cancellations}{The transformed assumptions and all particle-count cancellations}}


In scaled particle coordinates, the sum potential \(V_N\) has strong convexity \(Nm\) and gradient Lipschitz constant \(NL\). Put \(c=\sqrt N\), \(\gamma^\prime=c\gamma\), and \(h^\prime=h/c\). The original restrictions \(m+L\le\gamma^2\) and \(4\gamma Lh\le m\) become exactly \[Nm+NL\le(\gamma^\prime)^2,\qquad4\gamma^\prime(NL)h^\prime\le Nm.\] One proves positivity, curvature compatibility, friction, and step conditions together using \(c^2=N\), then proves the exact rate identity \[1-\frac{3(Nm)h^\prime}{4\gamma^\prime}=1-\frac{3mh}{4\gamma}.\] For \(D=d_\theta+Nd_x\), the shared bias simplifies exactly to \[\frac{NLh^\prime\sqrt{2D}}{Nm}=\frac{Lh}{m}\sqrt{\frac{2D}{N}}.\] The scalar proof checks the square-root quotient with nonnegative dimensions and the positive particle count. This cancellation explains why the first algorithm’s error contains \(d_x+d_\theta/N\), with no additional growing power of particle count.



Both curvature constants scale by \(N\), while friction and step scale by reciprocal square-root particle factors. The transformed rate cancels these factors exactly. The bias retains only thermal radius \(\sqrt{d_x+d_\theta/N}\), which explains the dependence on particle count.

\begin{leanbox}{Lean: \texttt{kiplmc1\_\allowbreak{}scaled\_\allowbreak{}conditions}, \texttt{kiplmc1\_\allowbreak{}scaled\_\allowbreak{}rate}, \texttt{kiplmc1\_\allowbreak{}scaled\_\allowbreak{}bias}}
\leansource{\texttt{kiplmc1\_\allowbreak{}scaled\_\allowbreak{}conditions} \quad KIPLMC/\allowbreak{}CanonicalExponential.lean, lines 42–60}
\begin{leancode}
\leanline{0}{theorem\ kiplmc1\_scaled\_conditions\ (N\ :\ ℕ)\ (hN\ :\ 0<N)\ (m\ :\ ℝ)\ (L\ γ\ h\ :\ ℝ≥0)}
\leanline{4}{(hm\ :\ 0<m)\ (hmL\ :\ m≤(L:ℝ))\ (hγ\ :\ 0<(γ:ℝ))}
\leanline{4}{(hfriction\ :\ m+(L:ℝ)≤(γ:ℝ)\textasciicircum{}2)\ (hh\ :\ 0<(h:ℝ))}
\leanline{4}{(hstep\ :\ 4*(γ:ℝ)*(L:ℝ)*(h:ℝ)≤m)\ :}
\leanline{4}{0<(N:ℝ)*m\ ∧\ (N:ℝ)*m≤((N:ℝ≥0)*L:ℝ≥0)\ ∧}
\leanline{4}{0<Real.sqrt\ N*(γ:ℝ)\ ∧}
\leanline{4}{(N:ℝ)*m+((N:ℝ≥0)*L:ℝ≥0)≤(Real.sqrt\ N*(γ:ℝ))\textasciicircum{}2\ ∧}
\leanline{4}{0<(h:ℝ)/Real.sqrt\ N\ ∧}
\leanline{4}{4*(Real.sqrt\ N*(γ:ℝ))*((N:ℝ≥0)*L:ℝ≥0)*((h:ℝ)/Real.sqrt\ N)≤(N:ℝ)*m}
\end{leancode}
\leansource{\texttt{kiplmc1\_\allowbreak{}scaled\_\allowbreak{}rate} \quad KIPLMC/\allowbreak{}CanonicalExponential.lean, lines 63–71}
\begin{leancode}
\leanline{0}{theorem\ kiplmc1\_scaled\_rate\ (N\ :\ ℕ)\ (hN\ :\ 0<N)\ (m\ γ\ h\ :\ ℝ)\ :}
\leanline{4}{1-3*((N:ℝ)*m)*(h/Real.sqrt\ N)/(4*(Real.sqrt\ N*γ))=1-3*m*h/(4*γ)}
\end{leancode}
\leansource{\texttt{kiplmc1\_\allowbreak{}scaled\_\allowbreak{}bias} \quad KIPLMC/\allowbreak{}CanonicalExponential.lean, lines 74–82}
\begin{leancode}
\leanline{0}{theorem\ kiplmc1\_scaled\_bias\ (N\ dθ\ dx\ :\ ℕ)\ (hN\ :\ 0<N)\ (m\ L\ h\ :\ ℝ)\ :}
\leanline{4}{((N:ℝ)*L)*(h/Real.sqrt\ N)*Real.sqrt\ (2*((dθ:ℝ)+(N:ℝ)*(dx:ℝ)))/((N:ℝ)*m)\ =}
\leanline{6}{L*h*(Real.sqrt\ 2*Real.sqrt\ (((dθ:ℝ)+(N:ℝ)*(dx:ℝ))/(N:ℝ)))/m}
\end{leancode}
\leanlabel{How the Lean proof follows the argument}
\texttt{kiplmc1\_\allowbreak{}scaled\_\allowbreak{}conditions} proves positivity of the real particle count and its square root as \texttt{hn} and \texttt{hc}; these allow multiplication of the original assumptions and cancellation in the transformed step. \texttt{kiplmc1\_\allowbreak{}scaled\_\allowbreak{}rate} handles zero friction separately before clearing nonzero denominators. \texttt{kiplmc1\_\allowbreak{}scaled\_\allowbreak{}bias} rewrites the square-root quotient, handles zero curvature separately, and then cancels the positive square-root particle factor. The last two scalar identities therefore have exactly the weak hypotheses stated in their own declarations, while the sampler application supplies positive curvature and friction.
\end{leanbox}

\section{\texorpdfstring{Derive finite transport moments for the output}{Derive finite transport moments for the output}}


The initial position law has the stated finite second moment. Strong convexity proves the corresponding moment of the normalized Gibbs target; their independent product therefore supplies a finite transport coupling. The independent product construction places this coupling on the same Gaussian/\allowbreak{}Brownian source as the chain. Its numerical and exact reference positions are \(L^{2}\) random variables, and exact reference stationarity has already been proved from the splitting limit. Their \(L^{2}\) discrepancy supplies a finite output transport coupling. The proof transfers this finiteness through the exact particle law identity. Thus the later real-valued Wasserstein and parameter RMS expressions are used only after their finiteness has been proved.



Finite initialization and target moments construct a finite-cost coupling. The shared constructed chain and stationary reference then give finite output transport on the source. This validates every later real-valued transport and parameter-error integral.

\begin{leanbox}{Lean: \texttt{kiplmc1GaussianPositionLaws\_\allowbreak{}wasserstein2\_\allowbreak{}ne\_\allowbreak{}top}}
\leansource{\texttt{kiplmc1GaussianPositionLaws\_\allowbreak{}wasserstein2\_\allowbreak{}ne\_\allowbreak{}top} \quad KIPLMC/\allowbreak{}CanonicalExponential.lean, lines 86–108}
\begin{leancode}
\leanline{0}{theorem\ kiplmc1GaussianPositionLaws\_wasserstein2\_ne\_top\ \{N\ dθ\ dx\ :\ ℕ\}\ (hN\ :\ 0<N)}
\leanline{4}{(U\ :\ IPLA.Joint\ dθ\ dx→ℝ)\ (m\ :\ ℝ)\ (hm\ :\ 0<m)\ (L\ γ\ :\ ℝ≥0)}
\leanline{4}{(hdiff\ :\ Differentiable\ ℝ\ U)\ (hsc\ :\ StrongConvexOn\ univ\ m\ U)}
\leanline{4}{(hLip\ :\ LipschitzWith\ L\ (gradient\ U))\ (h\ :\ ℝ≥0)}
\leanline{4}{(ν\ :\ Measure\ (State\ N\ dθ\ dx))\ [IsProbabilityMeasure\ ν]\ (hν\ :\ MemLp\ (id\ :\ State\ N\ dθ\ dx→State\ N\ dθ\ dx)\ 2\ ν)}
\leanline{4}{(n\ :\ ℕ)\ :}
\leanline{4}{ULA.wasserstein2\ ((kiplmc1GaussianLaws\ U\ γ\ h\ ((velocityTarget\ N\ dθ\ dx).prod\ ν)\ n).map\ Prod.snd)}
\leanline{6}{(IPLA.gibbsTarget\ U)≠⊤}
\end{leancode}
\leanlabel{How the Lean proof follows the argument}
\texttt{kiplmc1\allowbreak{}Gaussian\allowbreak{}Position\allowbreak{}Laws\_\allowbreak{}wasserstein2\_\allowbreak{}ne\_\allowbreak{}top} gives the formal counterpart of this step. The main intermediate facts are \texttt{hn}, \texttt{h\ensuremath{\pi}}, \texttt{hfinite}, \texttt{hglobal}, \texttt{ho} in that order. Named proof applications include \texttt{ho}; the adjacent proof shows their hypotheses and order of use.
\end{leanbox}


Earlier results used here: \cite[Section~3.17.1]{LangevinLean2026ULASC}; \cite[Section~3.17.3]{LangevinLean2026ULASC}; \cite[Section~3.21.2]{LangevinLean2026ULASC}; \cite[Section~3.26.15]{LangevinLean2026ULASC}; \cite[Section~3.26.3]{LangevinLean2026ULASC}; \cite[Section~3.1.1]{LangevinLean2026IPLA}; \cite[Section~3.1.3]{LangevinLean2026IPLA}; \cite[Section~3.1.4]{LangevinLean2026IPLA}; \cite[Section~3.3.2]{LangevinLean2026IPLA}; \cite[Section~3.5.1]{LangevinLean2026IPLA}; \cite[Section~3.57.17]{LangevinLean2026Kinetic}; \cite[Section~3.57.18]{LangevinLean2026Kinetic}.

\section{\texorpdfstring{Apply the shared kinetic theorem to the particle sampler}{Apply the shared kinetic theorem to the particle sampler}}


The shared theorem is applied to the sum potential, using its derived \(Nm\)-strong convexity and \(NL\)-smoothness. It proves, rather than assumes, the stationary Brownian reference, source independence, mild identities, local exponential-integrator error, scalar recurrence, and transport infimum. The argument supplies the transformed assumptions proved above, rewrites the output through the exact clock/\allowbreak{}Gaussian-law identity, counts the finite particle coordinates, and simplifies the rate and bias. The resulting position-law estimate is \[W_2(\nu_n,\pi_N)\le\sqrt2\left(1-\frac{3mh}{4\gamma}\right)^nW_2(\nu,\pi_N)+\frac{Lh\sqrt2}{m}\sqrt{\frac{d_\theta+Nd_x}{N}}.\] The initialization in this theorem is explicitly the independent thermal velocity law times the supplied position law; all algorithmic Gaussian correlations are those of the constructed noise measure.



The shared theorem is applied only after the exact noise law, whole-chain conjugacy, and transformed assumptions have been established. Rewriting its conclusion gives the physical-time particle sampling bound. Every term refers to the independent correlated-Gaussian algorithm.

\begin{leanbox}{Lean: \texttt{kiplmc1GaussianPositionLaws\_\allowbreak{}wasserstein\_\allowbreak{}bound}}
\leansource{\texttt{kiplmc1GaussianPositionLaws\_\allowbreak{}wasserstein\_\allowbreak{}bound} \quad KIPLMC/\allowbreak{}CanonicalExponential.lean, lines 111–143}
\begin{leancode}
\leanline{0}{theorem\ kiplmc1GaussianPositionLaws\_wasserstein\_bound\ \{N\ dθ\ dx\ :\ ℕ\}\ (hN\ :\ 0<N)}
\leanline{4}{(U\ :\ IPLA.Joint\ dθ\ dx→ℝ)\ (m\ :\ ℝ)\ (L\ γ\ :\ ℝ≥0)}
\leanline{4}{(hdiff\ :\ Differentiable\ ℝ\ U)\ (hsc\ :\ StrongConvexOn\ univ\ m\ U)}
\leanline{4}{(hLip\ :\ LipschitzWith\ L\ (gradient\ U))\ (hmL\ :\ m≤(L:ℝ))\ (hm\ :\ 0<m)}
\leanline{4}{(hγ\ :\ 0<(γ:ℝ))\ (hfriction\ :\ m+(L:ℝ)≤(γ:ℝ)\textasciicircum{}2)}
\leanline{4}{(h\ :\ ℝ≥0)\ (hh\ :\ 0<(h:ℝ))\ (hstep\ :\ 4*(γ:ℝ)*(L:ℝ)*(h:ℝ)≤m)}
\leanline{4}{(ν\ :\ Measure\ (State\ N\ dθ\ dx))\ [IsProbabilityMeasure\ ν]\ (hν\ :\ MemLp\ (id\ :\ State\ N\ dθ\ dx→State\ N\ dθ\ dx)\ 2\ ν)}
\leanline{4}{(n\ :\ ℕ)\ :}
\leanline{4}{ULA.wasserstein2Real\ ((kiplmc1GaussianLaws\ U\ γ\ h\ ((velocityTarget\ N\ dθ\ dx).prod\ ν)\ n).map\ Prod.snd)}
\leanline{6}{(IPLA.gibbsTarget\ U)\ ≤}
\leanline{6}{Real.sqrt\ 2*(1-3*m*(h:ℝ)/(4*(γ:ℝ)))\textasciicircum{}n*ULA.wasserstein2Real\ ν\ (IPLA.gibbsTarget\ U)+}
\leanline{6}{(L:ℝ)*(h:ℝ)*(Real.sqrt\ 2*Real.sqrt\ (((dθ:ℝ)+(N:ℝ)*(dx:ℝ))/(N:ℝ)))/m}
\end{leancode}
\leanlabel{How the Lean proof follows the argument}
\texttt{hdiff'}, \texttt{hsc'} and \texttt{h\allowbreak{}Lip'} establish the lifted sum-potential assumptions. \texttt{h\ensuremath{\pi}} and \texttt{hfinite} provide finite target and initial transport moments. The changed friction and step are \texttt{\ensuremath{\gamma}'} and \texttt{h'}; \texttt{kiplmc1\_\allowbreak{}scaled\_\allowbreak{}conditions} supplies all six transformed hypotheses. \texttt{hb} applies \texttt{Generic\allowbreak{}Canonical\allowbreak{}Coupling.\allowbreak{}klmc\_\allowbreak{}wasserstein\_\allowbreak{}bound}. The remaining rewrites use \texttt{kiplmc1\allowbreak{}Gaussian\allowbreak{}Position\allowbreak{}Laws\_\allowbreak{}eq\_\allowbreak{}canonical}, count the disjoint-union coordinates, and apply \texttt{kiplmc1\_\allowbreak{}scaled\_\allowbreak{}rate} and \texttt{kiplmc1\_\allowbreak{}scaled\_\allowbreak{}bias} to obtain the physical-time formula.
\end{leanbox}


Earlier results used here: \cite[Section~3.17.4]{LangevinLean2026ULASC}; \cite[Section~3.21.2]{LangevinLean2026ULASC}; \cite[Section~3.26.15]{LangevinLean2026ULASC}; \cite[Section~3.26.3]{LangevinLean2026ULASC}; \cite[Section~3.1.1]{LangevinLean2026IPLA}; \cite[Section~3.1.3]{LangevinLean2026IPLA}; \cite[Section~3.1.4]{LangevinLean2026IPLA}; \cite[Section~3.3.2]{LangevinLean2026IPLA}; \cite[Section~3.5.1]{LangevinLean2026IPLA}; \cite[Section~3.57.14]{LangevinLean2026Kinetic}; \cite[Section~3.57.19]{LangevinLean2026Kinetic}.

\section{\texorpdfstring{Project to the parameter and add the sharp equilibrium concentration}{Project to the parameter and add the sharp equilibrium concentration}}


The parameter projection is a continuous linear map with norm at most one. The powered Gibbs parameter marginal satisfies the proved sharp concentration bound \(\mathbb E_{\pi_N}\|\theta-\theta_\star\|^2\le d_\theta/(Nm)\). The transport-to-RMS inequality therefore gives \[\bigl(\mathbb E\|\theta_n-\theta_\star\|^2\bigr)^{1/2}\le\sqrt2\left(1-\frac{3mh}{4\gamma}\right)^nW_2(\nu,\pi_N)+\frac{Lh\sqrt2}{m}\sqrt{\frac{d_\theta+Nd_x}{N}}+\sqrt{\frac{d_\theta}{Nm}}.\] The argument invokes the output finiteness theorem, the sharp equilibrium concentration theorem, and the nonexpansive projection theorem. It rewrites the finite RMS as the displayed square root of the Bochner integral and substitutes the proved Gaussian position-law estimate. These are the initialization, discretization, and finite-particle concentration errors in the paper’s theorem.



The parameter projection is nonexpansive, so it cannot increase the transport discrepancy. Minkowski’s inequality adds the equilibrium parameter radius around the maximizer. This decomposes the parameter root-mean-square error into transient, discretization, and statistical contributions.

\begin{leanbox}{Lean: \texttt{kiplmc1\_\allowbreak{}parameter\_\allowbreak{}bound}}
\leansource{\texttt{kiplmc1\_\allowbreak{}parameter\_\allowbreak{}bound} \quad KIPLMC/\allowbreak{}CanonicalExponential.lean, lines 146–169}
\begin{leancode}
\leanline{0}{theorem\ kiplmc1\_parameter\_bound\ \{N\ dθ\ dx\ :\ ℕ\}\ (hN\ :\ 0<N)}
\leanline{4}{(U\ :\ IPLA.Joint\ dθ\ dx→ℝ)\ (m\ :\ ℝ)\ (L\ γ\ :\ ℝ≥0)}
\leanline{4}{(hdiff\ :\ Differentiable\ ℝ\ U)\ (hsc\ :\ StrongConvexOn\ univ\ m\ U)}
\leanline{4}{(hLip\ :\ LipschitzWith\ L\ (gradient\ U))\ (hmL\ :\ m≤(L:ℝ))\ (hm\ :\ 0<m)}
\leanline{4}{(hγ\ :\ 0<(γ:ℝ))\ (hfriction\ :\ m+(L:ℝ)≤(γ:ℝ)\textasciicircum{}2)}
\leanline{4}{(h\ :\ ℝ≥0)\ (hh\ :\ 0<(h:ℝ))\ (hstep\ :\ 4*(γ:ℝ)*(L:ℝ)*(h:ℝ)≤m)}
\leanline{4}{(ν\ :\ Measure\ (State\ N\ dθ\ dx))\ [IsProbabilityMeasure\ ν]\ (hν\ :\ MemLp\ (id\ :\ State\ N\ dθ\ dx→State\ N\ dθ\ dx)\ 2\ ν)}
\leanline{4}{(θStar\ :\ EuclideanSpace\ ℝ\ (Fin\ dθ))\ (hmax\ :\ IPLA.IsMarginalMaximizer\ U\ θStar)\ (n\ :\ ℕ)\ :}
\leanline{4}{Real.sqrt\ (∫\ z,\ ‖IPLA.parameterCLM\ z-θStar‖\textasciicircum{}2}
\leanline{6}{∂(kiplmc1GaussianLaws\ U\ γ\ h\ ((velocityTarget\ N\ dθ\ dx).prod\ ν)\ n).map\ Prod.snd)\ ≤}
\leanline{6}{Real.sqrt\ 2*(1-3*m*(h:ℝ)/(4*(γ:ℝ)))\textasciicircum{}n*ULA.wasserstein2Real\ ν\ (IPLA.gibbsTarget\ U)+}
\leanline{6}{(L:ℝ)*(h:ℝ)*(Real.sqrt\ 2*Real.sqrt\ (((dθ:ℝ)+(N:ℝ)*(dx:ℝ))/(N:ℝ)))/m+}
\leanline{6}{Real.sqrt\ ((dθ:ℝ)/((N:ℝ)*m))}
\end{leancode}
\leanlabel{How the Lean proof follows the argument}
\texttt{kiplmc1\_\allowbreak{}parameter\_\allowbreak{}bound} first derives finite output transport and invokes the exact parameter concentration bound. The nonexpansive parameter projection is passed to the projected root-mean-square transport inequality. Finiteness justifies rewriting the projected moment as its square-root integral. Substituting \texttt{kiplmc1\allowbreak{}Gaussian\allowbreak{}Position\allowbreak{}Laws\_\allowbreak{}wasserstein\_\allowbreak{}bound} then gives the transient, discretization and concentration terms in that order.
\end{leanbox}


Earlier results used here: \cite[Section~3.17.4]{LangevinLean2026ULASC}; \cite[Section~3.1.1]{LangevinLean2026IPLA}; \cite[Section~3.1.6]{LangevinLean2026IPLA}; \cite[Section~3.3.1]{LangevinLean2026IPLA}; \cite[Section~3.3.2]{LangevinLean2026IPLA}.

\section{\texorpdfstring{The paper’s independent-initialization constant}{The paper’s independent-initialization constant}}


The initial Wasserstein distance is bounded by every initial coupling cost. In particular, independent \(Z_0\sim\nu\) and \(\bar Z_\star\sim\pi_N\) give \[W_2(\nu,\pi_N)\le\bigl(\mathbb E\|Z_0-\bar Z_\star\|^2\bigr)^{1/2}.\] The argument constructs this independent product coupling, derives the integrability of its discrepancy from the two proved second moments, and rewrites its finite \(L^{2}\) cost as the square-root integral. It also proves that the stated step restrictions make the geometric rate nonnegative. Multiplying the initial coupling inequality by this nonnegative rate gives exactly the paper’s printed constant \(C_1\); the other two constants are unchanged. Thus the Wasserstein statement and the printed independent-initialization statement are both consequences for the same sampler.



The independent initial product coupling is one admissible plan, so its root-mean-square cost bounds the infimum defining Wasserstein distance. A nonnegative transient coefficient preserves that comparison. The printed initialization constant is therefore a corollary of the stronger Wasserstein estimate.

\begin{leanbox}{Lean: \texttt{independent\_\allowbreak{}initial\_\allowbreak{}wasserstein\_\allowbreak{}bound}, \texttt{kiplmc1\_\allowbreak{}rate\_\allowbreak{}nonneg}, \texttt{kiplmc1\_\allowbreak{}parameter\_\allowbreak{}bound\_\allowbreak{}independent}}
\leansource{\texttt{independent\_\allowbreak{}initial\_\allowbreak{}wasserstein\_\allowbreak{}bound} \quad KIPLMC/\allowbreak{}CanonicalExponential.lean, lines 173–191}
\begin{leancode}
\leanline{0}{theorem\ independent\_initial\_wasserstein\_bound}
\leanline{4}{\{E\ :\ Type*\}\ [NormedAddCommGroup\ E]\ [InnerProductSpace\ ℝ\ E]\ [CompleteSpace\ E]}
\leanline{4}{[MeasurableSpace\ E]\ [BorelSpace\ E]\ [SecondCountableTopology\ E]}
\leanline{4}{(ν\ π\ :\ Measure\ E)\ [IsProbabilityMeasure\ ν]\ [IsProbabilityMeasure\ π]}
\leanline{4}{(hν\ :\ MemLp\ (id\ :\ E→E)\ 2\ ν)\ (hπ\ :\ MemLp\ (id\ :\ E→E)\ 2\ π)\ :}
\leanline{4}{ULA.wasserstein2Real\ ν\ π\ ≤\ Real.sqrt\ (∫\ p\ :\ E×E,\ ‖p.1-p.2‖\textasciicircum{}2\ ∂ν.prod\ π)}
\end{leancode}
\leansource{\texttt{kiplmc1\_\allowbreak{}rate\_\allowbreak{}nonneg} \quad KIPLMC/\allowbreak{}CanonicalExponential.lean, lines 194–205}
\begin{leancode}
\leanline{0}{theorem\ kiplmc1\_rate\_nonneg\ (m\ :\ ℝ)\ (L\ γ\ :\ ℝ≥0)\ (hmL\ :\ m≤(L:ℝ))\ (hm\ :\ 0<m)}
\leanline{4}{(hγ\ :\ 0<(γ:ℝ))\ (hfriction\ :\ m+(L:ℝ)≤(γ:ℝ)\textasciicircum{}2)}
\leanline{4}{(h\ :\ ℝ≥0)\ (hstep\ :\ 4*(γ:ℝ)*(L:ℝ)*(h:ℝ)≤m)\ :}
\leanline{4}{0≤1-3*m*(h:ℝ)/(4*(γ:ℝ))}
\end{leancode}
\leansource{\texttt{kiplmc1\_\allowbreak{}parameter\_\allowbreak{}bound\_\allowbreak{}independent} \quad KIPLMC/\allowbreak{}CanonicalExponential.lean, lines 208–232}
\begin{leancode}
\leanline{0}{theorem\ kiplmc1\_parameter\_bound\_independent\ \{N\ dθ\ dx\ :\ ℕ\}\ (hN\ :\ 0<N)}
\leanline{4}{(U\ :\ IPLA.Joint\ dθ\ dx→ℝ)\ (m\ :\ ℝ)\ (L\ γ\ :\ ℝ≥0)}
\leanline{4}{(hdiff\ :\ Differentiable\ ℝ\ U)\ (hsc\ :\ StrongConvexOn\ univ\ m\ U)}
\leanline{4}{(hLip\ :\ LipschitzWith\ L\ (gradient\ U))\ (hmL\ :\ m≤(L:ℝ))\ (hm\ :\ 0<m)}
\leanline{4}{(hγ\ :\ 0<(γ:ℝ))\ (hfriction\ :\ m+(L:ℝ)≤(γ:ℝ)\textasciicircum{}2)}
\leanline{4}{(h\ :\ ℝ≥0)\ (hh\ :\ 0<(h:ℝ))\ (hstep\ :\ 4*(γ:ℝ)*(L:ℝ)*(h:ℝ)≤m)}
\leanline{4}{(ν\ :\ Measure\ (State\ N\ dθ\ dx))\ [IsProbabilityMeasure\ ν]\ (hν\ :\ MemLp\ (id\ :\ State\ N\ dθ\ dx→State\ N\ dθ\ dx)\ 2\ ν)}
\leanline{4}{(θStar\ :\ EuclideanSpace\ ℝ\ (Fin\ dθ))\ (hmax\ :\ IPLA.IsMarginalMaximizer\ U\ θStar)\ (n\ :\ ℕ)\ :}
\leanline{4}{Real.sqrt\ (∫\ z,\ ‖IPLA.parameterCLM\ z-θStar‖\textasciicircum{}2}
\leanline{6}{∂(kiplmc1GaussianLaws\ U\ γ\ h\ ((velocityTarget\ N\ dθ\ dx).prod\ ν)\ n).map\ Prod.snd)\ ≤}
\leanline{6}{Real.sqrt\ 2*(1-3*m*(h:ℝ)/(4*(γ:ℝ)))\textasciicircum{}n*}
\leanline{8}{Real.sqrt\ (∫\ p\ :\ State\ N\ dθ\ dx×State\ N\ dθ\ dx,\ ‖p.1-p.2‖\textasciicircum{}2\ ∂ν.prod\ (IPLA.gibbsTarget\ U))+}
\leanline{6}{(L:ℝ)*(h:ℝ)*(Real.sqrt\ 2*Real.sqrt\ (((dθ:ℝ)+(N:ℝ)*(dx:ℝ))/(N:ℝ)))/m+}
\leanline{6}{Real.sqrt\ ((dθ:ℝ)/((N:ℝ)*m))}
\end{leancode}
\leanlabel{How the Lean proof follows the argument}
\texttt{independent\_\allowbreak{}initial\_\allowbreak{}wasserstein\_\allowbreak{}bound} gives the formal counterpart of this step. The main intermediate facts are \texttt{hmem}, \texttt{hb}, \texttt{heq} in that order. Named proof applications include \texttt{heq}, \texttt{hb}; the adjacent proof shows their hypotheses and order of use.




\texttt{kiplmc1\_\allowbreak{}rate\_\allowbreak{}nonneg} gives the formal counterpart of this step. The main intermediate facts are \texttt{hcoeff}, \texttt{he}, \texttt{hp}, \texttt{hid} in that order. The final step is the algebraic normalization or ordered arithmetic displayed in the proof.




\texttt{kiplmc1\_\allowbreak{}parameter\_\allowbreak{}bound\_\allowbreak{}independent} gives the formal counterpart of this step. The main intermediate facts are \texttt{hn}, \texttt{h\ensuremath{\pi}}, \texttt{hw}, \texttt{hq}, \texttt{hc}, \texttt{hb} in that order. Named proof applications include \texttt{hb.\allowbreak{}trans}; the adjacent proof shows their hypotheses and order of use.
\end{leanbox}


Earlier results used here: \cite[Section~3.17.11]{LangevinLean2026ULASC}; \cite[Section~3.17.4]{LangevinLean2026ULASC}; \cite[Section~3.26.15]{LangevinLean2026ULASC}; \cite[Section~3.1.1]{LangevinLean2026IPLA}; \cite[Section~3.1.3]{LangevinLean2026IPLA}; \cite[Section~3.1.6]{LangevinLean2026IPLA}; \cite[Section~3.3.1]{LangevinLean2026IPLA}; \cite[Section~3.3.2]{LangevinLean2026IPLA}; \cite[Section~3.5.1]{LangevinLean2026IPLA}; \cite[Section~3.41.2]{LangevinLean2026Kinetic}.

\begin{thebibliography}{ULA-LSI}
\bibitem[ULA-SC]{LangevinLean2026ULASC}
LangevinLean Project. \emph{Unadjusted Langevin under Strong Convexity}. LL-01, 2026. Draft report, content revision d50c0f64815b. \url{https://odakyildiz.github.io/LangevinLean/docs/pdf/UnadjustedLangevin/FullProof.pdf}.
\bibitem[IPLA-SC]{LangevinLean2026IPLA}
LangevinLean Project. \emph{Interacting Particle Langevin under Strong Convexity}. LL-03, 2026. Draft report, content revision bfbb88fd1d79. \url{https://odakyildiz.github.io/LangevinLean/docs/pdf/IPLA/FullProof.pdf}.
\bibitem[Kinetic]{LangevinLean2026Kinetic}
LangevinLean Project. \emph{Kinetic Langevin Monte Carlo}. LL-06, 2026. Draft report, content revision f7f76041efde. \url{https://odakyildiz.github.io/LangevinLean/docs/pdf/KineticLangevin/FullProof.pdf}.
\end{thebibliography}

\end{document}
